Data-independent acquisition (DIA) is now the dominant mass-spectrometry strategy for proteomics at scale, and \texttt{DIA-NN}~\cite{demichev2020} is among the most widely adopted analysis engines across timsTOF~\cite{brunner2022,meier2020}, Astral, and Orbitrap platforms.
Recent benchmarks confirm \texttt{DIA-NN}'s competitiveness for single-cell DIA workflows~\cite{wang2025}, and recent releases have expanded support for plexDIA, timsTOF parallel accumulation-serial fragmentation (PASEF), and vendor-native raw-file reading.
Despite this analytical maturity, most \texttt{DIA-NN} users run the tool as a desktop application, which limits its scalability for large datasets. Deployments on HPC and cloud infrastructure remain improvised, relying on shell scripts, bespoke directory conventions, and hand-curated sample sheets. As a result, the same raw data analysed by two laboratories, or with two \texttt{DIA-NN} releases, can yield non-overlapping quantification tables that are difficult to reconcile.

Beyond per-study analysis, the value of public proteomics archives such as PRIDE Archive~\cite{perezriverol2025pride} and other ProteomeXchange partners~\cite{deutsch2026proteomexchange} lies as much in systematic reanalysis as in original deposition. In 2024 we published bigbio/quantms~\cite{dai2024quantms}, a Nextflow pipeline that standardised multi-engine DIA and DDA analyses behind a Sample and Data Relationship Format (SDRF)~\cite{dai2021sdrf} input on the nf-core framework~\cite{ewels2020nfcore}. quantms was, however, built around \texttt{DIA-NN}~1.8.1 with a fixed parameter set, and its DDA/DIA breadth made extending DIA-specific functionality difficult. New \texttt{DIA-NN} parameters conflicted with DDA options, versions were hard-coded, sample-level acquisition parameters were not properly propagated to per-run \texttt{DIA-NN} invocations, and raw-file conversion was mandatory. Iterative reanalysis at archive scale requires a workflow that combines the reproducibility of nf-core, the sensitivity of current \texttt{DIA-NN} releases, the metadata discipline of SDRF, and a parallelisation strategy capable of processing thousands of raw files.

Here, we present \quantmsdiann{}, an independent SDRF-driven Nextflow pipeline dedicated to \texttt{DIA-NN} that extends DIA-specific features beyond what the original bigbio/quantms pipeline could accommodate. While \quantmsdiann{} shares its SDRF-first design philosophy and parts of the configuration-generation tooling with quantms, its workflow, container set, and module library are developed and released independently, with container-pinned profiles for every \texttt{DIA-NN} version selectable at runtime under Docker or Singularity. Downstream, it exports a harmonised \qpx{} (Quantitative Proteomics eXchange) Parquet archive alongside the standard \texttt{DIA-NN} reports and MSstats input (Methods). We demonstrate the pipeline along three axes: ProteoBench~\cite{devreese2025proteobench} benchmarks across four DIA modalities and successive \texttt{DIA-NN} releases; independent reanalyses of public DIA deposits spanning bulk cell lines (NCI-60, ProCan), a single-cell oocyte plexDIA dataset, a spatial Deep Visual Proteomics cohort, and a phosphoproteomics cohort; and a pan-cohort of five public DIA cell-line reanalyses integrated under a single SDRF-driven invocation.
